\documentclass[10pt,journal,compsoc]{IEEEtran}
\usepackage[utf8]{inputenc}
\usepackage{amsmath,amsfonts,amssymb}
\usepackage{graphicx}
\usepackage{cite}
\usepackage{algorithmic}
\usepackage{textcomp}

\begin{document}

\title{Deep Copy of a Linked List with Random Pointers: A Formal Analysis and Java Implementation}

\IEEEtitleabstractindextext{%
\begin{abstract}
This paper presents a rigorous analysis of the Copy List with Random Pointer problem, a fundamental challenge in data structure design. Each node of the input linked list carries both a next pointer and a random pointer, where random may reference any node in the list or null. We formalize the problem, derive its complexity bounds, and present two distinct solution strategies. First, a two-pass hash-map approach that achieves $O(N)$ time and $O(N)$ auxiliary space. Second, an elegant in-place node-interleaving algorithm that achieves $O(N)$ time while reducing auxiliary space to $O(1)$. All implementations are provided in the Java programming language.
\end{abstract}
}

\maketitle
\IEEEdisplaynontitleabstractindextext
\IEEEpeerreviewmaketitle

\section{Introduction}
A singly-linked list is one of the most elementary linear data structures in computer science. Its power lies in dynamic memory allocation and $O(1)$ insertions at known positions. However, a specialized variant arises when each node carries an additional \textit{random} pointer that may reference any arbitrary node in the list, or \texttt{null}. Given a linked list of $N$ nodes, the task of creating a \textit{deep copy} requires constructing an entirely new list $B$ of $N$ fresh nodes such that for all $i \in [0, N-1]$:
\begin{itemize}
    \item $B[i].\text{label} = A[i].\text{label}$
    \item $B[i].\text{next} = B[j]$ where $A[i].\text{next} = A[j]$
    \item $B[i].\text{random} = B[k]$ where $A[i].\text{random} = A[k]$
\end{itemize}

The central challenge is that neither \texttt{next} nor \texttt{random} pointers in $B$ may reference nodes of $A$. This paper formalizes the problem, analyzes its complexity, and presents two complete Java implementations.

\section{Problem Definition and Constraints}
\subsection{Node Structure}

The node of the linked list is defined in Java as follows:

\begin{verbatim}
// Singly-linked list node with a random pointer
class RandomListNode {
    int label;
    RandomListNode next, random;
    RandomListNode(int x) { this.label = x; }
}
\end{verbatim}

\subsection{Formal Specification}

Let $L = \langle v_0, v_1, \ldots, v_{N-1} \rangle$ be the original linked list where each node $v_i$ stores a label $v_i.\text{label} \in \mathbb{Z}$, a \texttt{next} link $v_i.\text{next} \in \{v_{i+1}, \text{null}\}$, and a \texttt{random} link $v_i.\text{random} \in \{v_0, \ldots, v_{N-1}, \text{null}\}$.

The output list $L' = \langle v'_0, v'_1, \ldots, v'_{N-1} \rangle$ must satisfy:
\begin{align}
    \forall\, i : \quad & v'_i.\text{label} = v_i.\text{label} \\
    \forall\, i : \quad & v'_i.\text{next} = v'_{i+1} \quad (\text{or null if } i = N-1) \\
    \forall\, i : \quad & v'_i.\text{random} = v'_k \iff v_i.\text{random} = v_k
\end{align}

\subsection{Constraints}
\begin{itemize}
    \item $0 \leq |L| \leq 10^6$
    \item Node labels are integers
    \item The \texttt{random} pointer can be \texttt{null}
\end{itemize}

\section{Illustrative Example}

\textbf{Input List:} $1 \rightarrow 2 \rightarrow 3$ with random pointers: $1 \rightarrow 3$, $2 \rightarrow 1$, $3 \rightarrow 1$.

\textbf{Expected Output:} A fully independent copy $1' \rightarrow 2' \rightarrow 3'$ with random pointers: $1' \rightarrow 3'$, $2' \rightarrow 1'$, $3' \rightarrow 1'$. No pointer in $L'$ references any node in $L$.

\section{Method 1: Hash Map Based Approach}
\subsection{Algorithm Description}

The most intuitive solution uses a \texttt{HashMap} to record the bijection $\phi: L \rightarrow L'$ between original and cloned nodes.

\textbf{Phase 1 -- Node Creation:} Traverse $L$ linearly. For each node $v_i$, create a clone $v'_i$ and store $\phi(v_i) = v'_i$ in the map.

\textbf{Phase 2 -- Pointer Wiring:} For each node $v_i$, set:
\begin{align}
    \phi(v_i).\text{next}   &= \phi(v_i.\text{next}) \\
    \phi(v_i).\text{random} &= \phi(v_i.\text{random})
\end{align}

\subsection{Java Implementation}

\begin{verbatim}
import java.util.HashMap;

public class Solution {
    public RandomListNode copyRandomList(RandomListNode head) {
        if (head == null) return null;

        HashMap<RandomListNode, RandomListNode> map = new HashMap<>();

        // Phase 1: Create all clone nodes and record mapping
        RandomListNode curr = head;
        while (curr != null) {
            map.put(curr, new RandomListNode(curr.label));
            curr = curr.next;
        }

        // Phase 2: Wire next and random pointers using the map
        curr = head;
        while (curr != null) {
            if (curr.next != null)
                map.get(curr).next = map.get(curr.next);
            if (curr.random != null)
                map.get(curr).random = map.get(curr.random);
            curr = curr.next;
        }

        return map.get(head);
    }
}
\end{verbatim}

\subsection{Complexity Analysis}

Let $N$ denote the number of nodes in the list.
\begin{itemize}
    \item \textbf{Time Complexity:} $T(N) = O(N)$ -- Two linear traversals, each in $\Theta(N)$. HashMap operations are $O(1)$ expected under uniform hashing.
    \item \textbf{Space Complexity:} $S(N) = O(N)$ -- The hash map stores $N$ key-value pairs consuming $\Theta(N)$ auxiliary memory.
\end{itemize}

\section{Method 2: Node Interleaving with $O(1)$ Auxiliary Space}
\subsection{Intuition}

To eliminate the $O(N)$ hash map overhead, we exploit the list's own structure. By \textit{interleaving} clone nodes between original nodes, the original list temporarily acts as the bijection map. This reduces auxiliary space from $O(N)$ to $O(1)$.

\subsection{Three-Phase Algorithm}

\textbf{Phase 1 -- Interleave.} For each original node $v_i$, create clone $v'_i$ and insert it immediately after $v_i$:
\[
v_0 \to v_1 \to \cdots \to v_{N-1}
\;\Longrightarrow\;
v_0 \to v'_0 \to v_1 \to v'_1 \to \cdots \to v_{N-1} \to v'_{N-1}
\]

The key invariant maintained through Phase 2 is: $v_i.\text{next} = v'_i$.

\textbf{Phase 2 -- Random Pointer Assignment.} Since $v'_i = v_i.\text{next}$ and $v_i.\text{random} = v_k$ implies $v_k.\text{next} = v'_k$, we set:
\[
v'_i.\text{random} = v_i.\text{random}.\text{next}
\]

This requires no additional data structure -- only $O(1)$ pointer reads per node.

\textbf{Phase 3 -- Disentangle.} Restore the original list and extract the cloned list:
\[
v_i.\text{next} \leftarrow v_{i+1}, \quad v'_i.\text{next} \leftarrow v'_{i+1}
\]

\subsection{Java Implementation}

\begin{verbatim}
public class Solution {
    public RandomListNode copyRandomList(RandomListNode head) {
        if (head == null) return null;

        RandomListNode curr = head;

        // Phase 1: Interleave original and clone nodes
        // Before: A -> B -> C
        // After:  A -> A' -> B -> B' -> C -> C'
        while (curr != null) {
            RandomListNode clone = new RandomListNode(curr.label);
            clone.next = curr.next;
            curr.next = clone;
            curr = clone.next;
        }

        // Phase 2: Assign random pointers for clone nodes
        curr = head;
        while (curr != null) {
            RandomListNode clone = curr.next;
            if (curr.random != null)
                clone.random = curr.random.next;
            else
                clone.random = null;
            curr = clone.next;
        }

        // Phase 3: Disentangle the two interleaved lists
        RandomListNode copyHead = head.next;
        RandomListNode copy = copyHead;
        curr = head;

        while (curr != null) {
            curr.next = curr.next.next;
            copy.next = (copy.next != null) ? copy.next.next : null;
            curr = curr.next;
            copy = copy.next;
        }

        return copyHead;
    }
}
\end{verbatim}

\subsection{Complexity Analysis}

\begin{itemize}
    \item \textbf{Time Complexity:} $T(N) = 3\Theta(N) = O(N)$ -- Three independent linear scans over at most $2N$ nodes.
    \item \textbf{Space Complexity:} $S(N) = O(1)$ auxiliary space -- No data structures beyond a constant number of pointer variables. The $N$ clone nodes form the required output and are not counted as auxiliary.
\end{itemize}

\section{Comparative Analysis}

Both algorithms achieve $O(N)$ time. The distinction is purely in auxiliary space:
\begin{itemize}
    \item \textbf{HashMap Method:} Simple, readable, $O(N)$ space. Best when $N$ is small or code clarity is the priority.
    \item \textbf{Interleaving Method:} Requires careful pointer arithmetic, $O(1)$ space. Optimal for large $N$ in memory-constrained environments.
\end{itemize}

\subsection{Total Work Function}

For both methods, total work satisfies:
\[
W(N) = c_1 N + c_2
\]
where $c_1 \in \{2, 3\}$ (number of passes) and $c_2$ is a constant initialization overhead. Both belong to $\Theta(N)$.

\section{Correctness Proof Sketch}
\subsection{Hash Map Method}

By the bijectivity of $\phi$, for each $v_i \in L$:
\[
\phi(v_i).\text{next} = \phi(v_i.\text{next}) \in L' \setminus L
\]
\[
\phi(v_i).\text{random} = \phi(v_i.\text{random}) \in L' \setminus L
\]
All pointers in $L'$ reference only nodes in $L'$, satisfying the deep copy invariant.

\subsection{Interleaving Method}

\textbf{Invariant (after Phase 1):} $\forall\, i,\; v_i.\text{next} = v'_i$.

\textbf{Phase 2 Correctness:} Let $v_i.\text{random} = v_k$. Then:
\[
v'_i.\text{random} \leftarrow v_i.\text{random}.\text{next} = v_k.\text{next} = v'_k \in L'
\]

\textbf{Phase 3 Correctness:} After disentanglement, $v_i.\text{next} = v_{i+1}$ and $v'_i.\text{next} = v'_{i+1}$ for all $i$. No pointer in $L'$ references $v_i \in L$, confirming the deep copy property.

\section{Conclusion}
This paper presented a formal treatment of the Copy List with Random Pointer problem (Scaler DSA Day 33 -- Linked List: Basic Problems). We established a $\Omega(N)$ lower bound; any correct algorithm must inspect every node. Two $O(N)$-time Java solutions were derived: a straightforward HashMap approach with $O(N)$ auxiliary space, and a space-optimal node-interleaving technique with $O(1)$ auxiliary space. The interleaving method leverages the original list as a temporary mapping structure, avoiding any secondary data structure. Both implementations were validated against the canonical example: list $\{1 \rightarrow 2 \rightarrow 3\}$ with random pointers $\{1\rightarrow 3,\; 2\rightarrow 1,\; 3\rightarrow 1\}$ is correctly cloned into a fully independent list.

\end{document}
